\section{Introduction}
\label{sec:intro}
\todo{Motivation: graphs encode relationships (social, financial, molecular, biomedical) that are the sensitive part.}
\todo{Gap: membership inference on graphs is well studied; reconstruction of the graph itself from each kind of release is not. State precisely what is missing.}
\todo{Contributions (3 bullets): attack(s), a unified threat-model taxonomy, empirical results across GNNs and synthetic graph generators.}

\section{Background and Related Work}
\label{sec:related}
\paragraph{Graph learning and graph synthetic data.}
A graph is $G=(V,E,X)$ with node set $V$, edge set $E$, and node features $X\in\R^{|V|\times d}$. GNNs~\citep{kipf2017semi} learn by message passing over $E$. \todo{Synthetic graph generators: classical (ER, SBM, configuration model), deep (GraphRNN-style, diffusion-based), DP variants.}

\paragraph{Privacy attacks.}
Membership inference~\citep{shokri2017membership} asks whether a record was used; reconstruction asks for the record itself. \todo{Related work on link stealing, graph reconstruction from embeddings/gradients, model inversion, attribute inference on graphs.}

\begin{table}[t]
  \centering
  \caption{Working taxonomy: what the adversary tries to reconstruct (rows) vs.\ what is released (columns). Fill in which cells are covered by prior work and which by us.}
  \label{tab:taxonomy}
  \small
  \begin{tabular}{@{}lccc@{}}
    \toprule
    Target & GNN outputs & GNN weights/grads & Synthetic graph \\
    \midrule
    Links (edges)        & \todo{} & \todo{} & \todo{} \\
    Node attributes      & \todo{} & \todo{} & \todo{} \\
    Node labels          & \todo{} & \todo{} & \todo{} \\
    Graph-level (whole)  & \todo{} & \todo{} & \todo{} \\
    \bottomrule
  \end{tabular}
\end{table}

\section{Problem Setup and Threat Model}
\label{sec:threat}
\todo{Define: private graph $G$, release $\mathcal{R}(G)$, adversary knowledge (auxiliary graph? partial edges? model access: black-box/white-box?), adversary goal (reconstruct $E$, $X$, labels, or $G$ up to isomorphism).}
\todo{Metrics: edge AUC / precision@$k$ / F1, attribute accuracy or RMSE, graph edit distance or degree/clustering statistics for graph-level.}

\section{Attack Design}
\label{sec:attack}
\subsection{Attack A (working name)}
\todo{Intuition first, then algorithm. Use an \texttt{algorithm} block if it is procedural.}
\subsection{Attack B (working name)}
\todo{}

\section{Experimental Setup}
\label{sec:setup}
\todo{Datasets (Cora/Citeseer/PubMed, OGB, molecular sets for graph-level), target models (GCN, GraphSAGE, GAT, GIN), synthetic graph generators, baselines, seeds, compute.}

\section{Results}
\label{sec:results}
\todo{Main table + one figure. Report mean $\pm$ std over seeds.}
% \begin{figure}[t]
%   \centering
%   \includegraphics[width=\linewidth]{figs/main_result}
%   \caption{}
%   \label{fig:main}
% \end{figure}

\section{Discussion, Defenses, and Limitations}
\label{sec:discussion}
\todo{Does DP (edge-level vs node-level) stop these attacks? What do defenses cost?}
\todo{Limitations and ethical considerations.}
